\begin{tikzpicture}[
  >=Latex,
  x=1.25cm,
  y=1.25cm,
  line cap=round,
  line join=round,
  font=\small,
  frente/.style={draw=blue!65!black, line width=1.25pt},
  propagacion/.style={->, orange!85!black, line width=1.35pt},
  dimension/.style={<->, black!75, line width=0.9pt}
]

\coordinate (C) at (0,0);

% Frentes de onda en el corte radial del medio.
\foreach \radio in {1.35,1.95,2.55,3.15} {
  \draw[frente] (C) circle (\radio);
}

% Cavidad y radios característicos.
\fill[blue!3] (C) circle (0.86);
\draw[black!70, line width=1.4pt] (C) circle (0.86);
\node[font=\small\bfseries] at (0,0.28) {cavidad};
\fill[black!70] (C) circle (1.5pt);
\draw[dimension] (C) -- (0.78,-0.36)
  node[pos=0.66, below=1pt, fill=white, inner sep=1.5pt] {$R$};
\draw[dimension] (C) -- (-2.12,-2.12)
  node[pos=0.72, below left, fill=white, inner sep=1.5pt] {$r$};

% Flechas radiales: propagación hacia afuera.
\foreach \angulo in {30,90,150,270,330} {
  \draw[propagacion]
    ({1.05*cos(\angulo)},{1.05*sin(\angulo)}) --
    ({1.62*cos(\angulo)},{1.62*sin(\angulo)});
}

% Identificación de las ondas y del material.
\node[align=center, text=blue!55!black, fill=white, inner sep=2pt]
  at (-4.35,2.10) {frentes de onda\\longitudinal saliente};
\draw[blue!55!black, thin] (-3.42,1.80) -- (-2.05,2.05);
\node[align=center, fill=white, inner sep=2pt] at (4.15,-2.05)
  {medio elástico\\infinito};

\end{tikzpicture}